% Complete Section 8.3, original PDF 109-110 / printed 99-100.
% block: lower-order-terms-heading page:109 type:section
\subsection{{低阶项}}\label{sec:lower-order-terms}

% block: magnetostatic-lower-order-models pages:109,110 type:ordinary-paragraph
{静磁学中的一个基本模型问题, 是求 (向量代理) 2-形式 $B$, 使得}
\[
\operatorname{curl}\operatorname{curl}B=f,\qquad\operatorname{div}B=0\quad\text{in }\Omega
\]
{其中给定的 $f$ 属于旋度的像, 并对 $B$ 施加边界条件. 可以通过求解 Hodge Laplace 方程得到解}
\[
\operatorname{curl}\operatorname{curl}B-\operatorname{grad}\operatorname{div}B=f.
\]
{对时谐问题, 此方程中出现额外的零阶项, 得到扰动的 Hodge Laplace 方程}
\[
\operatorname{curl}\operatorname{curl}B-\operatorname{grad}\operatorname{div}B-m^2B=f,
\]
{其中 $m$ 是波数. 在磁流体力学中, 还有速度向量场 $v$ 引起的额外一阶对流项, 因而静态情形下有}
\[
\operatorname{curl}\operatorname{curl}B-\operatorname{grad}\operatorname{div}B-\operatorname{curl}(v\times B)=f.
\]

% block: general-lower-order-perturbations page:110 type:ordinary-paragraph
{出现这样的低阶项时, 作用于 $B$ 的微分算子无法实现为任何 Hilbert 复形的抽象 Hodge Laplacian. 然而, 将所研究的用于得到 Hodge Laplacian 可行数值方法的混合有限元方法, 推广到这类带低阶项的问题, 并不困难. 不太清楚的是这些离散化的收敛分析, 因为前面的分析不能直接应用于它们. 在 \BookCite{14} 中, 我们研究了 FEEC 混合方法对如下形式微分算子的收敛性}
\[
Lu=[(d+l_1)(\delta+l_2)+(\delta+l_3)(d+l_4)+l_5],
\]
{其中包含 Hodge Laplace 算子 $d\delta+\delta d$ 最自然的低阶扰动. 这里 $l_i$ 是以光滑系数场作乘子的算子. 例如, 时谐问题中的零阶项包含在 $l_5$ 中, 而磁流体力学问题中的一阶扰动由 $l_2B=-v\times B$ 给出 (2-形式与向量场的缩并).}

% block: perturbed-problem-stability-and-rates page:110 type:ordinary-paragraph
{\BookCite{14} 的分析相当复杂, 这里只简要概述主要结果. 算子 $L$ (包括其边界条件) 当然可能有核, 就像通常的 Helmholtz 算子 $-\Delta+l$ 在 $l$ 为负时可能有核一样. 假定并非如此, 即 $0$ 不是 $L$ 的特征值, 因而连续问题适定. 则为 Hodge Laplacian 导出的有限元离散化, 即 \eqref{eq:four-galerkin-form-pairs} 中的全部四种选择, 对扰动算子 $L$ 也稳定, 至少对充分细的剖分如此. 因此, Theorem~\ref{thm:mixed-basic-error} 陈述的基本误差估计成立. 意外出现在试图建立 Theorem~\ref{thm:mixed-improved-errors} 的改进误差估计时, 这些估计为 $\sigma$, $d\sigma$, $u$ 和 $du$ 各量建立最优 $L^2$ 收敛阶. 结果表明, 对有限元空间的某些选择以及某些低阶项 $(l_1,\ldots,l_5)$, $\sigma$ 或 $d\sigma$ 或两者可能损失一阶收敛阶. 结果总结于 Table~\ref{tab:perturbed-hodge-error-rates}, 应将其与 Table~\ref{tab:galerkin-form-approximation-rates} (所有 $l_i$ 均为零的情形) 比较. \BookCite{14} 在理论上建立了所有这些阶数, 并通过计算验证了其尖锐性, 细节请参见该文.}

% block: perturbed-hodge-error-rates-table page:110 type:table
\begin{table}[htbp]\centering\small\setlength{\tabcolsep}{3pt}
\caption{{Hodge Laplacian 的 FEEC 解的 $L^2$ 误差阶.}}\label{tab:perturbed-hodge-error-rates}
\begin{tabular}{cc|cccc}\hline
$V_h^{k-1}$&$V_h^k$&$\sigma$&$d\sigma$&$u$&$du$\\\hline
$\mathcal P_r\Lambda^{k-1}$&$\mathcal P_r^-\Lambda^k$&$\begin{cases}r+1&\text{if }l_2=l_5=0,\\r&\text{otherwise}.\end{cases}$&$r$&$r$&$r$\\\noalign{\vskip4pt}
$\mathcal P_r^-\Lambda^{k-1}$&$\mathcal P_r^-\Lambda^k$&$r$&$r$&$r$&$r$\\\noalign{\vskip4pt}
$\mathcal P_r\Lambda^{k-1}$&$\mathcal P_{r-1}\Lambda^k$&$\begin{cases}r+1&\text{if }l_2=l_3=l_5=0,\\r&\text{otherwise}.\end{cases}$&$\begin{cases}r&\text{if }l_3=0,\\r-1&\text{otherwise}.\end{cases}$&$r$&$r-1$\\\noalign{\vskip4pt}
$\mathcal P_r^-\Lambda^{k-1}$&$\mathcal P_{r-1}\Lambda^k$&$r$&$\begin{cases}r&\text{if }l_3=0,\\r-1&\text{otherwise}.\end{cases}$&$r$&$r-1$\\\hline
\end{tabular}
\end{table}
